YPF ha implementado cartelería digital en sus estaciones, lo que permite desplegar 
comunicación comercial dinámica: publicidades, promociones y mensajes segmentados 
por momento del día, perfil de cliente o contexto de compra. Este módulo analiza si 
esa infraestructura tiene un impacto medible sobre las ventas de combustibles y 
productos de tienda, e identifica qué tipos de mensajes y momentos de exposición 
resultan más efectivos.

\subsection{Objetivo}

Las preguntas centrales que guían el análisis son:

\begin{itemize}
    \item ¿La cartelería digital aumenta las ventas en la estación?
    \item ¿Qué tipo de mensajes generan mayor respuesta: publicidades de marca 
    o promociones de producto?
    \item ¿Importa el momento del día en que se muestra el mensaje? ¿Una 
    publicidad de café a las 8:00 AM aumenta la probabilidad de comprar 
    un café en la tienda?
    \item ¿La exposición a publicidad en la cartelería aumenta la interacción 
    con la aplicación móvil?
    \item ¿Qué retorno económico tiene la inversión en cartelería digital?
\end{itemize}

\subsection{Diseño econométrico}

El principal desafío metodológico es que las estaciones que cuentan con cartelería 
digital no son necesariamente comparables con las que no la tienen: pueden diferir 
en tamaño, ubicación o perfil de clientes. Además, el contenido que se muestra en 
cada estación no es aleatorio, sino que responde a decisiones comerciales que pueden 
estar correlacionadas con las ventas. Comparar directamente estaciones con y sin 
cartelería, o períodos con y sin determinada publicidad, podría entonces reflejar esas 
diferencias previas en lugar del efecto real de la comunicación.

La estrategia para sortear este problema aprovecha dos fuentes de variación. La 
primera es la rotación del contenido publicitario dentro de una misma estación a lo 
largo del tiempo: si una estación alterna entre distintos mensajes en distintos 
momentos del día, es posible comparar el comportamiento de los clientes expuestos 
a cada tipo de contenido controlando por las características fijas de la estación. La 
segunda fuente de variación es la instalación escalonada de la cartelería en distintas 
estaciones, que permite aplicar una lógica similar a la del módulo de Tiendas Full: 
comparar cada estación consigo misma antes y después de la instalación, usando como 
referencia estaciones aún no equipadas en ese período.

\subsubsection*{Ejercicio 1: Efecto de la cartelería sobre la interacción con 
la aplicación y sus consecuencias sobre ventas}

\textbf{Pregunta:} ¿La exposición a publicidad en la cartelería aumenta el uso de 
la aplicación móvil? ¿Y ese mayor uso se traduce en mayores ventas?

\textbf{Lógica del ejercicio:} Este ejercicio aborda una cadena causal en dos 
etapas. En la primera, se estima si la exposición a la cartelería aumenta la 
interacción del cliente con la aplicación —medida por tiempo de uso, apertura o 
transacciones realizadas a través de ella—. En la segunda, se usa esa variación 
inducida por la publicidad como palanca para estimar el efecto del uso de la 
aplicación sobre las ventas. Este enfoque, conocido como variables instrumentales, 
permite aislar la parte del uso de la aplicación que se explica por la exposición 
publicitaria —y que por lo tanto puede considerarse exógena— para estimar su 
impacto causal sobre el comportamiento de compra.

\textbf{Datos necesarios:} Datos de interacción con la aplicación a nivel de 
usuario, estación y hora; registro de contenidos en cartelería digital; 
transacciones de combustibles y tienda.


\subsubsection*{Ejercicio 2: Impacto del contenido publicitario sobre ventas 
en tienda}

\textbf{Pregunta:} ¿La exposición a un mensaje específico en la cartelería digital 
aumenta la probabilidad de comprar el producto promocionado?

\textbf{Lógica del ejercicio:} Para cada producto o categoría que aparece en la 
cartelería, se estima la probabilidad de compra en función de si el mensaje 
correspondiente estaba activo en el momento de la visita, controlando por efectos 
fijos de estación, día y hora, así como por el flujo general de clientes. Esto permite 
comparar el comportamiento de clientes en condiciones similares, diferenciados 
únicamente por el mensaje al que estuvieron expuestos. La granularidad horaria de 
los datos es clave para este ejercicio: permite capturar, por ejemplo, si una 
publicidad de café activa durante la mañana tiene un efecto distinto al de la misma 
publicidad activada al mediodía.

\textbf{Datos necesarios:} Registro de contenidos mostrados en cartelería digital 
por estación y franja horaria; transacciones de tienda con fecha, hora y producto; 
datos de tráfico de clientes por estación y período.



\subsection{Datos requeridos}

El análisis requiere datos desagregados a nivel de estación con frecuencia horaria 
o diaria. En particular, es necesario contar con el registro de contenidos mostrados 
en la cartelería —tipo de mensaje, producto o categoría promocionada, franja horaria 
de activación— que deberá construirse junto con el equipo responsable de la 
comunicación comercial. Esta categorización es un insumo indispensable para poder 
distinguir el efecto de distintos tipos de mensajes. Las fuentes transaccionales 
requeridas se detallan en la sección~\ref{sec:datos}.

